% ============================================================
\section{Kết luận}
% ============================================================
Bài báo đề xuất \textbf{StatLM}, một pipeline encoder--decoder cho bài toán phân
loại văn bản đa nhãn zero-shot trong môi trường thế giới mở cho tiếng Việt. Bằng
cách giao việc giảm nhiễu và cô đặc thông tin cho encoder thống kê (TextRank trên
sentence embedding, KeyBERT với embedding bất đối xứng) và chỉ huy động LLM một lần ở
decoder, StatLM cân bằng chất lượng, chi phí và tốc độ: so với baseline X-MLClass-Gemini dùng
LLM toàn phần, nó nâng tỉ lệ nhãn thô lên 76,9\%, hạ nhiễu xuống 3,8\%, phủ 19/20
chuyên mục, và giảm token cùng thời gian khoảng $30\times$ và $4\times$.

\textbf{Hạn chế.} Chất lượng phụ thuộc dây chuyền vào hai tầng encoder; chi phí xây
đồ thị tăng $O(n^2)$ với văn bản rất dài; đầu ra LLM còn dao động nhẹ giữa các lần
chạy; và tham chiếu keyphrase là LLM-proxy chứ chưa phải ground-truth do người gán.
Chuỗi xử lý cũng là vòng hở: một chủ đề phụ chỉ nằm trong vài câu có độ trung tâm
thấp có thể bị bỏ mà không được khôi phục---một rủi ro \emph{được cơ chế dự đoán}
dưới benchmark một nhãn vàng ($|L_i|=1$), mà để định lượng cần một tập con đa nhãn do
người gán.

\textbf{Hướng phát triển.} Tích hợp thông tin vị trí câu và cấu trúc văn bản vào đồ
thị; mở rộng sang các miền khác và xây dựng bộ dữ liệu keyphrase tiếng Việt gán nhãn
thủ công; và bổ sung cơ chế dự phòng dựa trên độ tin cậy, đưa lại toàn văn vào khâu
trích keyphrase khi độ phủ bản tóm tắt hoặc độ tin cậy gán nhãn thấp, khép một phần
vòng lặp hiện hở.

